\chapter{Execution Architecture: From C Processes to Python Runtime Semantics}

A Python course can begin with Python syntax, but that approach leaves the student with a dangerous illusion: that programming is mainly the act of arranging visible language constructs. It is not. Programming is the act of describing computations that must eventually be represented, loaded, scheduled, executed, interrupted, resumed, and stored inside a real machine. A student who only learns the visible syntax of Python may become comfortable writing small scripts, but will often remain unable to explain why variables behave differently from C variables, why function calls create frames, why lists store references rather than raw values, why exceptions can unwind several calls at once, why generators can resume after yielding, or why asynchronous programs can handle many waiting tasks without using many CPU cores.

This chapter therefore starts below Python. It begins with the native execution model, using C as the reference architecture. This is not a detour. It is the baseline. C shows the traditional path clearly: source code becomes object code, object code becomes an executable, the operating system loader maps that executable into memory, the process receives a stack and heap, and the CPU executes machine instructions from the text segment. Once this model is understood, Python can be explained honestly: not as magic, not as ``just interpreted'', and not as a loose collection of convenient syntax rules, but as a runtime-hosted language executed inside a native interpreter process.

The analogy is simple. One may learn to move a car in a parking lot without understanding roads, intersections, lanes, bridges, roundabouts, traffic lights, and right-of-way rules. But that person is not yet a competent driver. In the same way, one may learn to print strings, write loops, and define functions without understanding processes, memory, frames, objects, bytecode, and runtime scheduling. That person can operate the surface of the language, but cannot yet understand the execution environment in which real programs live.

The first section of the chapter establishes the native process model. It explains what a program is, what a process is, what the operating system kernel does, how compilation transforms C source code into executable form, how loading creates an active virtual address space, and how memory is divided into text, data, BSS, heap, and stack regions. This gives the student a concrete mental model for compiled execution: where instructions live, where variables live, how function calls use stack frames, why heap allocation must be managed, and why the operating system remains in control even when the program is native machine code.

The second section introduces process virtual machines and interpreted runtimes. This is where the Python execution model becomes visible. When a Python file is run, the operating system does not load that file as a native executable. It loads the CPython executable. The Python file is input to a runtime. CPython parses it, compiles it to bytecode, creates code objects and frames, allocates Python objects, and executes virtual instructions through an interpreter loop. This section also compares CPython with the JVM, V8, PHP, and Bash, because the word \emph{interpreted} is too vague to be useful unless we separate source interpretation, bytecode interpretation, opcode execution, JIT compilation, host-provided event loops, and command dispatch.

The third section studies modern procedural extensions: callbacks, anonymous functions, closures, generators, coroutines, and asynchronous execution. These constructs are placed here because they only make full sense after the runtime model is clear. A closure requires preserved lexical state beyond an ordinary stack frame. A generator requires suspended execution state. A coroutine requires controlled resumption. An asynchronous task requires the runtime or event loop to remember what should continue after an external event becomes ready. These are not just elegant language features. They are consequences of treating code, frames, continuations, and waiting operations as runtime-managed structures.

The chapter is therefore laid out as a bridge. It moves from native execution to runtime-hosted execution, and then from ordinary procedures to runtime-managed control flow. This order is deliberate. Python's variable model cannot be explained well before the object model. The object model cannot be explained well before the heap and references. Python frames cannot be explained well before native stack frames. Python bytecode cannot be explained well before machine code and instruction streams. Async execution cannot be explained well before scheduling, blocking, and suspended state. The chapter builds these dependencies in the order in which they become intelligible.

The purpose is not to turn a Python book into an operating systems book or a compiler construction book. The purpose is to give Python the execution context it deserves. A serious programming language is not only its syntax. It is a complete execution architecture. Python becomes much easier to understand when seen as a managed runtime layered above the native process model inherited from systems such as C.

This approach also protects the student from misleading simplifications. Python variables are not C boxes. Python assignment is not memory copying. Python functions are not only named source blocks. Python exceptions are not only error messages. Python generators are not lists. Python concurrency is not automatically parallelism. Each of these facts becomes clear only when the underlying architecture is visible.

For that reason, this chapter is not preliminary decoration. It is the foundation of the course. A student who understands this chapter will read the rest of Python differently. Syntax will no longer be memorized as isolated rules. It will be interpreted as surface notation for deeper mechanisms: objects, bindings, frames, bytecode, namespaces, managed memory, and runtime-controlled execution.